\documentclass[11pt,a4paper]{article}
\usepackage[T1]{fontenc}
\usepackage[utf8]{inputenc}
\usepackage{lmodern}
\usepackage[margin=27mm,headheight=14pt]{geometry}
\usepackage{amsmath,amssymb,amsthm,mathtools,mathrsfs}
\usepackage{microtype,booktabs,array,longtable,aliascnt}
\usepackage{enumitem}
\usepackage{xcolor}
\usepackage{fancyhdr}
\usepackage{hyperref}
\usepackage[nameinlink,noabbrev]{cleveref}
\hypersetup{colorlinks=true,linkcolor=blue!40!black,citecolor=blue!40!black,
 urlcolor=blue!40!black,pdftitle={Global Infinitesimal Surcomplex Dynamics},
 pdfauthor={AI-assisted research manuscript},pdfsubject={Hahn-coherent dynamics and period invariants}}
\pagestyle{fancy}
\fancyhf{}
\fancyhead[L]{\small Global infinitesimal surcomplex dynamics}
\fancyhead[R]{\small Research manuscript}
\fancyfoot[C]{\thepage}
\setlength{\parskip}{4pt}
\setlength{\parindent}{1.3em}
\setlist{itemsep=2pt,topsep=4pt}
\emergencystretch=1.5em
\numberwithin{equation}{section}
\newtheorem{theorem}{Theorem}[section]
\newaliascnt{proposition}{theorem}
\newtheorem{proposition}[proposition]{Proposition}
\aliascntresetthe{proposition}
\newaliascnt{lemma}{theorem}
\newtheorem{lemma}[lemma]{Lemma}
\aliascntresetthe{lemma}
\newaliascnt{corollary}{theorem}
\newtheorem{corollary}[corollary]{Corollary}
\aliascntresetthe{corollary}
\theoremstyle{definition}
\newaliascnt{definition}{theorem}
\newtheorem{definition}[definition]{Definition}
\aliascntresetthe{definition}
\newaliascnt{example}{theorem}
\newtheorem{example}[example]{Example}
\aliascntresetthe{example}
\theoremstyle{remark}
\newaliascnt{remark}{theorem}
\newtheorem{remark}[remark]{Remark}
\aliascntresetthe{remark}
\crefname{proposition}{proposition}{propositions}
\Crefname{proposition}{Proposition}{Propositions}
\crefname{lemma}{lemma}{lemmas}
\Crefname{lemma}{Lemma}{Lemmas}
\crefname{corollary}{corollary}{corollaries}
\Crefname{corollary}{Corollary}{Corollaries}
\crefname{definition}{definition}{definitions}
\Crefname{definition}{Definition}{Definitions}
\crefname{example}{example}{examples}
\Crefname{example}{Example}{Examples}
\crefname{remark}{remark}{remarks}
\Crefname{remark}{Remark}{Remarks}
\newcommand{\C}{\mathbb C}
\newcommand{\Q}{\mathbb Q}
\newcommand{\R}{\mathbb R}
\newcommand{\N}{\mathbb N}
\newcommand{\Z}{\mathbb Z}
\newcommand{\No}{\mathbf{No}}
\newcommand{\A}{\mathcal A}
\newcommand{\G}{\mathcal G}
\newcommand{\Oh}{\mathcal O}
\newcommand{\Mh}{\mathcal M}
\newcommand{\mm}{\mathfrak m}
\newcommand{\DD}{\mathscr D}
\newcommand{\eps}{\varepsilon}
\newcommand{\id}{\mathrm{id}}
\newcommand{\supp}{\operatorname{supp}}
\newcommand{\Res}{\operatorname{Res}}
\newcommand{\st}{\operatorname{st}}
\newcommand{\Per}{\operatorname{Per}}
\newcommand{\val}{\operatorname{v}}
\newcommand{\Span}{\operatorname{span}}
\newcommand{\dd}{\,dz}
\newcommand{\Fexp}[1]{\exp(#1\partial_z)z}
\newcommand{\pos}{_{>0}}
\newcommand{\nonneg}{_{\geq0}}
\newcommand{\Oe}[1]{O(\eps^{#1})}
\newcommand{\snapshot}{39f2be6667ade51bca2b45daa47e289d69c09764}

\title{\textbf{Global Infinitesimal Surcomplex Dynamics}\\[5pt]
\large Complete period invariants, difference equations,\\
and moduli beyond every power of one infinitesimal}
\author{AI-assisted research manuscript}
\date{21 September 2026}

\begin{document}
\maketitle
\begin{abstract}
We study near-identity holomorphic dynamics with coefficients in a set-sized
Hahn field $K=\C((t^\Gamma))$, allowing an arbitrary ordered value group.
All Hahn coefficient functions are holomorphic on one fixed ordinary domain
$D\subset\C$.  This common-domain requirement makes global analytic questions
meaningful without importing an inappropriate notion of convergence on
$\No[i]$.

The principal result classifies maps with a fixed nowhere-vanishing leading
infinitesimal displacement: two such maps are conjugate by a globally defined
Hahn-coherent near-identity change of variable precisely when the periods of
their canonical time forms agree.  A basepoint normalization makes the
conjugacy unique.  On a plane with $r$ punctures, the resulting moduli space is
exactly $\mm_K^r$.  We also identify the cokernel of the discrete difference
operator with $K^r$ and give an explicit Bernoulli-operator solution whenever
the period obstructions vanish.

In higher rank these results produce a family of pairwise nonconjugate maps
that are indistinguishable modulo every power of a chosen infinitesimal.
Their exact period invariants are nonzero at a smaller scale.  A further
theorem constructs all Hahn coefficients of a slow-time flow on any specified
ordinary flow chart, without successive shrinking of that chart.

The exponential--logarithm correspondence for summable automorphisms is
established background, not a novelty claim.  The proposed contributions are
the common-domain global classification, its explicit obstruction theory,
and the higher-rank invisible-moduli construction.  Full proofs are supplied;
novelty is assessed by a targeted, not exhaustive, literature search.  No
named published conjecture is claimed solved, and the manuscript is not
independently refereed or formally verified.
\end{abstract}

\newpage
{\fontsize{10}{11}\selectfont\setlength{\parskip}{0pt}\tableofcontents}
\newpage

\section{The problem and the main conclusions}
\label{sec:introduction}

An infinitesimal displacement can be small at every ordinary point of a
complex domain and nevertheless carry a global obstruction.  The distinction
is not about the magnitude of the displacement.  It is about whether locally
constructed coordinates agree after continuation around an ordinary loop.
This observation suggests studying surcomplex dynamics through differential
forms rather than through sequences of iterates.

Fix an ordered abelian group $\Gamma$ which is a set and a connected ordinary
complex domain $D$.  The function algebra used throughout is
\begin{equation}\label{eq:algebra-intro}
 \A_\Gamma(D)=\Oh(D)((t^\Gamma)).
\end{equation}
Its elements are Hahn series in which \emph{every} coefficient is holomorphic
on the same $D$.  The variable is $z$; the Hahn monomials are constants.
An infinitesimal map is
\[
 F(z)=z+h(z),\qquad h\in\A_\Gamma(D),\quad \val(h)>0.
\]
Although $F$ need not be an ordinary complex-valued self-map of $D$, it is a
well-defined map on its infinitesimal halo.  Its inverse is defined in the
same category.  Composition is Taylor substitution justified by support
finiteness, not by taking limits of partial sums.

For such a map there is a canonical positive Hahn vector field $W_F$ with
\begin{equation}\label{eq:log-intro}
 F=\exp(W_F\partial_z)z.
\end{equation}
The general algebraic mechanism behind this formula is known.  We prove the
particular version needed here, including its global-domain and support
properties, in \cref{sec:logarithm}.

Suppose that for some $\delta>0$ and some nowhere-zero $a\in\Oh(D)$,
\begin{equation}\label{eq:stratum-intro}
 F(z)=z+t^\delta a(z)+r(z),\qquad \val(r)>\delta.
\end{equation}
Then $W_F=t^\delta(a+u)$ with $\val(u)>0$.  Thus $W_F$ is invertible in
$\A_\Gamma(D)$ and the \emph{time form}
\[
 \Omega_F=\frac{dz}{W_F}
\]
is a globally defined Hahn-coherent holomorphic one-form on $D$.
The adjective ``holomorphic'' here refers to all its ordinary coefficient
functions; its valuation may be negative.

\subsection{A classification by periods}

The main equivalence, proved in \cref{thm:classification}, is
\begin{equation}\label{eq:main-intro}
 \begin{split}
 &H\circ F_1=F_2\circ H\quad\text{for some }H=z+\text{positive Hahn terms}
 \\
 &\hspace{12mm}\Longleftrightarrow\quad
 \oint_\gamma\Omega_{F_1}=\oint_\gamma\Omega_{F_2}
 \quad\text{for every ordinary closed curve }\gamma\subset D.
 \end{split}
\end{equation}
Both maps must have the same data $(\delta,a)$ in
\eqref{eq:stratum-intro}.  No shrinking of $D$ is permitted in this
statement.  All the contour integrals are coefficientwise ordinary complex
integrals.  For every $z_*\in D$, there is exactly one conjugacy in
\eqref{eq:main-intro} satisfying $H(z_*)=z_*$.

For
\[
 D=\C\setminus\{a_1,\ldots,a_r\},
\]
only $r$ periods are required.  After subtracting the leading time form and
normalizing by $t^\delta$, they can be prescribed arbitrarily in
$\mm_K^r$; see \cref{thm:moduli}.  Thus the classification is not merely
an obstruction: it gives an explicit parametrization and representatives
of all the conjugacy classes in the specified stratum.

\subsection{Discrete cohomology and a global Abel equation}

Writing $\Delta_F\psi=\psi\circ F-\psi$, we show that
\begin{equation}\label{eq:exact-intro}
 0\longrightarrow K\longrightarrow\A_\Gamma(D)
 \xrightarrow{\Delta_F}\A_\Gamma(D)
 \xrightarrow{\mathcal P_F}K^r\longrightarrow0
\end{equation}
is exact on a finitely punctured plane, where
\[
 (\mathcal P_Fg)_j=\frac1{2\pi i}\oint_{\gamma_j}\frac{g(z)}{W_F(z)}\,dz.
\]
In particular, the global Abel equation
\[
 \psi\circ F=\psi+1
\]
has a solution in the common-domain algebra exactly when every period of
$\Omega_F$ vanishes.  A solution is an Abel function, not automatically a
globally injective coordinate.  This distinction will be retained throughout.

\subsection{A phenomenon genuinely involving more than one scale}

Suppose $\eta>n\delta$ for every positive ordinary integer $n$.  Put
$\eps=t^\delta$ and choose constants $c_1,\ldots,c_r$ whose valuations are
larger than every $n\delta$.  The fields and maps
\begin{equation}\label{eq:flat-intro}
 W_{\mathbf c}(z)=
 \frac{\eps}{1+\eps\displaystyle\sum_{j=1}^r c_j/(z-a_j)},
 \qquad F_{\mathbf c}=\exp(W_{\mathbf c}\partial_z)z
\end{equation}
satisfy
\[
 F_{\mathbf c}\equiv z+\eps\pmod{\eps^n\A_\Gamma(D)_{\geq0}}
 \quad\text{for every }n\geq1.
\]
Nevertheless,
\begin{equation}\label{eq:flat-time-intro}
 \Omega_{F_{\mathbf c}}=\eps^{-1}dz+
       \sum_{j=1}^r\frac{c_j\,dz}{z-a_j}
\end{equation}
shows that distinct parameter vectors give nonconjugate maps.  Their Abel
obstructions are exactly $(c_1,\ldots,c_r)$.  Thus a complete collection of
one-scale jets can omit exact global dynamical information.
This is proved in \cref{thm:flat}, including the valuation estimate and a
concrete realization inside $\No[i]$.

\section{Context, provenance, and the boundary of the claims}
\label{sec:context}

\subsection{Relation to the supplied repository}

The repository was consulted at the snapshot
\begin{center}
\small\ttfamily 39f2be6667ade51bca2b45daa47e289d69c09764.
\end{center}
Its catalogue and the relevant analysis and contour-calculus documentation
separate common-domain Hahn functions from radius-free coefficient germs,
and coefficientwise contour integration from fine-topological convergence
\cite{repo-catalogue,repo-analysis,repo-contours}.  We use that distinction,
but do not assume the global geometric assertions of the repository's other
research drafts.  In particular, the argument below does not depend on an
unproved surcomplex Nullstellensatz, a contour representation for a
multivariate residue, or a general global Picard theorem.

The review was of the catalogue and the relevant analysis/contour material,
not a line-by-line audit of every manuscript in the repository.  The subject
selected here is the global dynamics of near-identity maps and their time
forms, rather than another treatment of local Taylor calculus, polynomial
root clusters, or divisor existence.

\subsection{Established mathematics used here}

The exponential correspondence between suitably summable derivations and
automorphisms of generalized-series algebras was developed in general form
by Bagayoko, Krapp, Kuhlmann, Panazzolo, and Serra
\cite{bkkps}.  Accordingly, the mere existence of a formal logarithm in
\eqref{eq:log-intro} is \emph{not} proposed as new.  Our proof specializes
that mechanism to holomorphic coefficient functions on one fixed domain and
makes the finite support dependencies explicit.

Classical iterative logarithms also have a substantial theory.  The work of
Aschenbrenner and Bergweiler \cite{ab} treats Julia's equation and
transcendence properties of iterative logarithms.  Our parameter-Hahn
construction is a different category: no convergence in an ordinary
nonzero parameter is asserted.  It does not remove the analytic embedding
obstructions for ordinary parabolic germs, nor classify their analytic
sectorial invariants.

Modified differential equations in backward error analysis provide another
close antecedent.  Hairer's lectures \cite{hairer} explain their formal
construction and warn that the ordinary step-size series generally diverges.
The Euler coefficients displayed in \cref{sec:examples} are known examples,
not new coefficients.  Their role here is to exhibit an exact global period
obstruction in the Hahn category.  The slow-time result in
\cref{sec:slowtime} concerns coefficientwise holomorphic existence on a
specified ordinary flow chart, not analytic convergence of the modified
series in a real step size.

\subsection{What is proposed as a contribution}

The main proposed result is the complete \emph{global}, \emph{common-domain},
\emph{arbitrary-rank} classification in \cref{thm:classification}, together
with the explicit moduli realization in \cref{thm:moduli}.  The discrete
cohomology calculation, the higher-rank invisible-moduli family, and the
common-chart slow-time construction are complementary results.  The proofs
combine known algebraic and analytic ingredients, but the precise package
of statements was not located in the targeted search.

That assessment is deliberately narrower than a priority claim.  An
unlocated result may exist in literature on formal flows, deformation of
one-forms, or generalized series.  The manuscript therefore offers the
stated theorems and their proofs as candidate new results, rather than
claiming to settle a named published open problem.  None of the central
claims is presented as an independently checked theorem of Lean.

\section{The common-domain Hahn calculus}
\label{sec:hahn}

\subsection{Fields, coefficients, and summability}

Let $\Gamma$ be a totally ordered abelian group which is a set.  For a
commutative $\C$-algebra $R$, write
\[
 R((t^\Gamma))=
 \left\{\sum_{\gamma\in S}r_\gamma t^\gamma:
              S\subset\Gamma\text{ is well ordered}\right\}.
\]
Zero coefficients are discarded from the support.  An element has one
coefficient at each exponent, and no infinite ordinary sum is hidden in a
single coefficient.  Addition and multiplication are the Hahn operations.
When $R$ is an integral domain, the valuation of a nonzero series is its
least support exponent; set $\val(0)=+\infty$.

We will use
\[
 K=\C((t^\Gamma)),\qquad
 \A=\Oh(D)((t^\Gamma)),\qquad
 \mathcal B=\Mh(D)((t^\Gamma)),
\]
where $D$ is nonempty, open, and connected, and $\Mh(D)$ is the field of
ordinary meromorphic functions on $D$.  Thus $K$ and $\mathcal B$ are
fields, but $\A$ is generally not a field.  Constants embed as $K\subset\A$.
Write
\[
 \A_{\geq0}=\{f:\val(f)\geq0\},\quad
 \A_{>0}=\{f:\val(f)>0\},\quad
 \mm_K=\{c\in K:\val(c)>0\}.
\]
These sets include zero under our convention for its valuation.

\begin{definition}[Strong summability]\label{def:summability}
A set-indexed family $(f_i)_{i\in I}$ in $R((t^\Gamma))$ is strongly
summable if the union of its supports is well ordered and, for every
$\gamma\in\Gamma$, only finitely many $i$ have a nonzero coefficient at
$\gamma$.  Its sum is then defined coefficientwise by finite sums in $R$.
A map is strongly linear if it is linear and preserves these summable
families and their sums.
\end{definition}

The support condition, not a metric estimate, is the basis of every infinite
operator expression below.  In particular, knowing that an operator raises
valuation by $\delta>0$ does not by itself justify a sequential argument
which silently assumes that $n\delta$ is cofinal in $\Gamma$.

\subsection{A support lemma with a finite-word proof}

For $S\subset\Gamma_{>0}$, let $S^*$ be the set of all finite sums of
members of $S$, including the empty sum $0$.

\begin{lemma}[Positive-support lemma]\label{lem:neumann}
If $S$ is well ordered and positive, then $S^*$ is well ordered, and every
$\gamma\in\Gamma$ is the sum of only finitely many ordered words with
letters in $S$.  If $T\subset\Gamma$ is well ordered, then $T+S^*$ is
well ordered and every $\gamma$ has only finitely many decompositions
$\gamma=\tau+\mu$ with $\tau\in T$, $\mu\in S^*$.
\end{lemma}

\begin{proof}
We give the usual finite-word argument to emphasize that no Archimedean
hypothesis is involved.  A word $u$ is embedded in a word $w$ if the letters
of $u$ can be matched, in their order, to a subsequence of the letters of
$w$, with each matched letter no larger than its partner.

First, every infinite sequence of words over a well-ordered alphabet has
an earlier word embedded in a later word.  Otherwise choose a bad sequence
$w_0,w_1,\ldots$, at each position choosing a word of least possible length
among those allowing a bad continuation after the chosen prefix.  No word
is empty.  Write $w_n=v_n a_n$, with last letter $a_n$.  The sequence of
last letters has an infinite nondecreasing subsequence
$a_{n_0}\leq a_{n_1}\leq\cdots$: successively choose a minimum in each
remaining infinite tail.  The sequence
\[
 w_0,\ldots,w_{n_0-1},v_{n_0},v_{n_1},v_{n_2},\ldots
\]
cannot be bad, by minimality of the length of $w_{n_0}$.  An embedding
among the original prefix words, or from one of them into a $v_{n_j}$,
contradicts the original badness.  An embedding $v_{n_j}\hookrightarrow
v_{n_k}$ extends to $w_{n_j}\hookrightarrow w_{n_k}$ by matching the last
letters.  This is again a contradiction.

When the letters are positive elements of $\Gamma$, an embedding implies
that the sum of the earlier word is at most the sum of the later word.
The inequality is strict unless the words are identical: an additional
letter contributes positively, and a strictly increased matched letter
also increases the sum.  Consequently an infinite strictly descending
sequence of word sums is impossible.  Also an infinite sequence of distinct
words with the same sum is impossible.  This proves both assertions about
$S^*$.

Finally, any infinite sequence in a well-ordered set has an infinite
nondecreasing subsequence.  Applying this successively to the coordinates
of pairs in $T\times S^*$ shows that the product order admits no infinite
sequence with no comparable earlier/later pair.  A strictly descending
sequence of pair sums would have that property.  Infinitely many distinct
pairs with a fixed sum would have it as well, since a comparable pair with
the same sum must have equal coordinates.  This proves the last assertions.
\end{proof}

\begin{corollary}[Finite dependency]\label{cor:dependency}
Suppose an expression is built from a seed series with support $T$ by
finitely many operations at each stage, and every additional perturbative
factor contributes an exponent from a fixed well-ordered positive set
$S$.  A sum over all finite such perturbative words is strongly summable,
provided each fixed word contributes only finitely many terms.  Its support
is contained in $T+S^*$.
\end{corollary}

\begin{proof}
For a prescribed output exponent there are finitely many seed/word
possibilities by \cref{lem:neumann}.  Each of these has only finitely many
internal contributions.  The union of output supports is contained in the
well-ordered set $T+S^*$.
\end{proof}

We use this corollary also for derivatives, which do not introduce new Hahn
exponents.  The derivative order at a fixed output exponent is finite,
although no common finite derivative order is asserted for the whole series.

\subsection{Differentiation, units, and primitives}

Define
\[
 \partial_z\left(\sum f_\gamma(z)t^\gamma\right)
       =\sum f_\gamma'(z)t^\gamma.
\]
This is a strongly $K$-linear derivation, and its kernel is $K$, because an
ordinary holomorphic function with zero derivative is constant on connected
$D$.  The same kernel statement holds in $\mathcal B$.

\begin{lemma}[A sufficient unit criterion]\label{lem:unit}
If $f=t^\lambda(a+u)$ with $a\in\Oh(D)$ nowhere zero and
$u\in\A_{>0}$, then $f$ is invertible in $\A$, with
\[
 f^{-1}=t^{-\lambda}a^{-1}\sum_{n\geq0}(-u/a)^n.
\]
Every coefficient of the inverse is holomorphic on the original $D$.
\end{lemma}
\begin{proof}
The displayed family is summable by \cref{lem:neumann}.  Multiplication by
$1+u/a$ telescopes coefficientwise to $1$.  Powers of $a^{-1}$ are
holomorphic on $D$, so no new singularity is introduced there.
\end{proof}

We make no converse assertion about arbitrary units of $\A$; only the
stated criterion is needed.  Every nonzero series is invertible in
$\mathcal B$, since its leading nonzero coefficient is invertible as a
meromorphic function.

\begin{lemma}[Support-preserving primitives]\label{lem:primitive}
For $g\in\A$, a primitive $P\in\A$ with $P'=g$ exists if and only if
\[
 \oint_\gamma g(z)\,dz=0
\]
for every ordinary closed piecewise $C^1$ curve in $D$, where the integral
is taken coefficientwise.  If it exists, the primitive with $P(z_*)=0$
is unique and introduces no new support exponents.
\end{lemma}
\begin{proof}
Necessity holds coefficientwise by the ordinary fundamental theorem for
line integrals.  Conversely define each coefficient primitive by path
integration from $z_*$ to $z$.  Vanishing periods make this independent of
the path.  On an ordinary disk the resulting function is an ordinary
holomorphic primitive, hence it is holomorphic on all of $D$.  Use zero for
a coefficient that was zero.  The resulting support is a subset of
$\supp(g)$, so the collection is a Hahn series.  Any two primitives differ
by a member of $K$, and the normalization removes that constant.
\end{proof}

\subsection{Taylor substitution and actual halo values}

For $h\in\A_{>0}$, put $F=z+h$ and define
\begin{equation}\label{eq:substitution}
 T_Fg=g\circ F:=\sum_{n\geq0}\frac{h^n}{n!}\,g^{(n)}.
\end{equation}
With $S=\supp(h)$ and $T=\supp(g)$, this is strongly summable, with support
in $T+S^*$.  It is strongly $K$-linear.  Leibniz's formula proves that it
is an algebra homomorphism, and the usual Taylor chain rule gives
\[
 (g\circ F)'=(g'\circ F)F'.
\]
These identities may be checked at each exponent, where there are only
finitely many terms; there is no unjustified rearrangement of ordinary
infinite sums.  The same argument proves associativity of substitution.
Notice the operator-order convention
\begin{equation}\label{eq:antiorder}
 T_FT_G=T_{G\circ F}.
\end{equation}

For an actual point of the halo
\[
 D^\#=\{z_0+\xi:z_0\in D,\ \xi\in\mm_K\}\subset K,
\]
define
\begin{equation}\label{eq:halo-eval}
 g(z_0+\xi)=\sum_{\gamma,n}
      \frac{g_\gamma^{(n)}(z_0)}{n!}\,t^\gamma\xi^n.
\end{equation}
The same support lemma applies.  No common lower bound on the ordinary
Taylor radii of infinitely many functions is required to justify the Hahn
sum; each coefficient is differentiated only finitely often for a fixed
output exponent.  Equation \eqref{eq:halo-eval} respects sums, products, and
substitution.  It sends $F=z+h$ to a map of $D^\#$ into itself, because
$\st(F(z_0+\xi))=z_0$.

Evaluation is faithful for this coefficient-function category.  If
$g\ne0$, choose an ordinary point where its leading nonzero coefficient
function is nonzero; then $g$ has nonzero value at that point.  Thus the
identities below are identities of genuine halo functions as well as of
Hahn coefficient families.

\subsection{Why this gives surcomplex results}

When $\Gamma$ is an additive subgroup of $\No$, the standard Conway normal
form identification sends $t^\gamma$ to $\omega^{-\gamma}$.  A Hahn
series with complex coefficients is sent to its real and imaginary surreal
normal forms.  Their supports are sets, so this gives a field embedding
\[
 \C((t^\Gamma))\hookrightarrow\No[i].
\]
This uses the classical normal-form arithmetic of the surreal numbers
\cite{gonshor}; it does not require a set containing all surreal numbers.
The positive-support sums used in this paper remain set-indexed after the
embedding.  One can also enlarge $\Gamma$ to contain the supports of a
specified set of data and then work entirely in that workspace.

The particular rank-two example in \cref{sec:flat} needs only
$\Gamma=\Q+\Q\omega$.  The resulting theorems therefore already describe
explicit surcomplex maps.  They are not statements about arbitrary class
functions on the entire surcomplex plane, and $\partial_z$ is not an
intrinsic derivation of the scalar field $\No$.

\section{Canonical logarithms and exact infinitesimal suspension}
\label{sec:logarithm}

Let
\[
 \G_+(D)=\{z+h:h\in\A_{>0}\}.
\]
We will prove that this is a group, rather than assuming its invertibility
properties.  For $F=z+h$ put $\Delta=T_F-I$.

\begin{theorem}[Common-domain logarithm--exponential correspondence]
\label{thm:logexp}
For every $F\in\G_+(D)$, the strongly summable operator
\begin{equation}\label{eq:operatorlog}
 L_F=\log T_F=\sum_{n\geq1}\frac{(-1)^{n+1}}{n}\Delta^n
\end{equation}
has the form $L_F=W_F\partial_z$ for a unique $W_F\in\A_{>0}$.
Conversely, every $W\in\A_{>0}$ defines
\begin{equation}\label{eq:operatorexp}
 T=\exp(W\partial_z)=\sum_{n\geq0}\frac{(W\partial_z)^n}{n!},
 \qquad F=Tz\in\G_+(D),
\end{equation}
and $T=T_F$.  The two constructions are inverse bijections.  They preserve
the original coefficient domain $D$.  If $h\ne0$, then
\[
 \val(W_F)=\val(h),\qquad
 \text{the leading coefficient of }W_F\text{ equals that of }h.
\]
The inverse of $F$ is $\exp(-W_F\partial_z)z$.
\end{theorem}

\begin{proof}
\emph{Summability.}
From Taylor substitution,
\[
 \Delta g=\sum_{k\geq1}\frac{h^k}{k!}g^{(k)}.
\]
Every term of $\Delta^n g$ contains at least $n$ perturbative factors from
$h$ or its derivatives.  At a fixed Hahn exponent, the allowed words in
$\supp(h)$ and the seed exponent of $g$ are finite in number, by
\cref{cor:dependency}.  Consequently the sum over $n$ in
\eqref{eq:operatorlog} is strongly summable.  Its coefficients, on an
ordinary input function, are finite-order differential operators with
holomorphic coefficients on $D$.  The same reasoning applies to the
exponential of $W\partial_z$, using $\supp(W)$.

\emph{The logarithm is a derivation.}
Introduce an ordinary indeterminate $u$ and the binomial expression
\[
 U_u=\sum_{n\geq0}\binom un\Delta^n.
\]
For every fixed output exponent on fixed inputs, this is a polynomial in
$u$, not an infinite power series in $u$.  When $u=N$ is a nonnegative
integer, $U_N=T_F^N$ is multiplicative.  Hence the coefficient of every
Hahn exponent in
\[
 U_u(fg)-(U_uf)(U_ug)
\]
is a polynomial which vanishes at all nonnegative integers.  It is zero.
Differentiating this polynomial identity at $u=0$, and using
\[
 \left.\frac{d}{du}\binom un\right|_{u=0}
   =\frac{(-1)^{n+1}}n\quad(n\geq1),
\]
proves the Leibniz identity for $L_F$.

\emph{The derivation is a vector field.}
Let $W_F=L_Fz$.  On an ordinary holomorphic input, the coefficient of any
fixed exponent in $L_F$ is a finite-order differential operator $P$ and
is a derivation.  It annihilates constants.  At any $z_0\in D$, apply
Leibniz to powers of $z-z_0$.  The value of $P$ on $(z-z_0)^j$ at $z_0$
is zero for $j\ne1$, and the value for $j=1$ is the coefficient of $W_F$
at $z_0$.  These polynomials determine a finite jet, so
$P f(z_0)=Pz(z_0)f'(z_0)$ for every holomorphic $f$.  Varying $z_0$ gives
$L_F=W_F\partial_z$ on ordinary coefficients.  Strong $K$-linearity
extends this identity to $\A$.

This finite-jet step is important: no claim about arbitrary abstract
$\C$-derivations of a function algebra is being smuggled into the proof.
Only the explicitly constructed differential operators are involved.

\emph{Exponentiating a vector field.}
Repeated Leibniz gives
\[
 (W\partial_z)^n(fg)=\sum_{j=0}^n\binom nj
    ((W\partial_z)^j f)((W\partial_z)^{n-j}g).
\]
After a summable rearrangement this proves that
$\exp(W\partial_z)$ is multiplicative and fixes $K$.
Set $F=\exp(W\partial_z)z$.  It equals $z+$ positive terms.
On polynomials in $z$, the exponential homomorphism equals substitution by
$F$.  At each Hahn exponent, both operations are finite-order differential
operators on an ordinary holomorphic input.  Equality on polynomial jets
at every point therefore implies equality on all holomorphic inputs, and
then on $\A$ by strong linearity.

\emph{Inverse identities and the leading term.}
The formal scalar identities
$\exp(\log(1+X))=1+X$ and $\log(\exp X)=X$ hold here coefficientwise:
at each exponent their compositions use only finitely many perturbative
words.  Thus the constructions are inverse.  The first term of $L_Fz$
is $h$, and every remaining term contains at least two factors from
$h$ and its derivatives.  If $\val(h)=\delta>0$, the remainder has
valuation at least $2\delta>\delta$.  The leading-term assertion follows.
Finally $\exp(W\partial_z)\exp(-W\partial_z)=I$, which gives the
stated inverse map with the composition convention
\eqref{eq:antiorder}.
\end{proof}

\begin{remark}[What ``exact suspension'' means]
The theorem supplies a strongly summable vector field whose time-one
substitution is exactly $F$.  It does not claim that an ordinary complex
specialization of the Hahn parameter has a convergent iterative logarithm.
Nor does it construct a continuous trajectory for the fine topology on
all of $\No[i]$.  The analytic extension to an ordinary slow-time chart
will require a separate construction in \cref{sec:slowtime}.
\end{remark}

\begin{corollary}[Julia equation]\label{cor:julia}
For every $F\in\G_+(D)$,
\begin{equation}\label{eq:julia}
 W_F(F(z))=F'(z)W_F(z).
\end{equation}
\end{corollary}
\begin{proof}
$T_F$ commutes with its logarithm.  Apply
$T_F(W_F\partial_z)=(W_F\partial_z)T_F$ to $z$.
\end{proof}

\begin{proposition}[Conjugacy and vector fields]\label{prop:intertwine}
For $F_1,F_2,H\in\G_+(D)$, the following are equivalent:
\[
 H\circ F_1=F_2\circ H,
 \qquad W_{F_2}(H)=H'W_{F_1}.
\]
\end{proposition}
\begin{proof}
The conjugacy relation is $T_{F_1}T_H=T_HT_{F_2}$.  Conjugating the
summable logarithm gives
$(W_{F_1}\partial_z)T_H=T_H(W_{F_2}\partial_z)$.
Evaluation on $z$ gives the displayed vector-field identity.  Conversely
that identity and the chain rule give the operator intertwining on all
of $\A$.  Exponentiation recovers the conjugacy.
\end{proof}

\section{Time forms and their global periods}
\label{sec:timeforms}

Assume for the moment that $W=W_F$ has a nowhere-zero leading coefficient.
By \cref{lem:unit}, the time form
\[
 \Omega_F=W^{-1}dz
\]
belongs to $\A\,dz$.  Its significance is elementary: in a local primitive
coordinate $P$ with $P'=1/W$, the vector field $W\partial_z$ becomes a
unit translation vector field.  The issue is whether such local primitives
fit together globally and whether one can compare the primitives of two
maps without leaving the common coefficient domain.

For $H=z+h\in\G_+(D)$ define the pullback of a one-form by
\[
 H^*(g\,dz)=(g\circ H)H'\,dz.
\]
All operations are those already justified in \cref{sec:hahn}.

\begin{lemma}[Explicit exactness of an infinitesimal pullback]
\label{lem:homotopy}
For $g\in\A$ and $h\in\A_{>0}$, the series
\begin{equation}\label{eq:Q}
 Q_h(g)=\sum_{n\geq1}\frac{h^n}{n!}g^{(n-1)}
\end{equation}
is strongly summable and satisfies
\begin{equation}\label{eq:pullback-exact}
 (z+h)^*(g\,dz)-g\,dz=dQ_h(g).
\end{equation}
Consequently every coefficientwise period is unchanged by a near-identity
pullback.
\end{lemma}
\begin{proof}
For $S=\supp(h)$, the support in \eqref{eq:Q} is contained in
$\supp(g)+(S^*\setminus\{0\})$, and each coefficient has finitely many
contributions.  Termwise differentiation is therefore legitimate.  It gives
\begin{align*}
 (Q_h(g))'
 &=\sum_{n\geq1}\frac{h^n}{n!}g^{(n)}
   +h'\sum_{n\geq1}\frac{h^{n-1}}{(n-1)!}g^{(n-1)}\\
 &=g(z+h)-g(z)+h'g(z+h)\\
 &=g(z+h)(1+h')-g(z).
\end{align*}
This is \eqref{eq:pullback-exact}.  Exact coefficient forms have zero
ordinary periods by \cref{lem:primitive}.
\end{proof}

This lemma is a support-controlled version of the familiar fact that an
infinitesimal change of coordinate acts trivially on one-dimensional
holomorphic de Rham cohomology.  The proof is included as infrastructure;
the homotopy principle itself is not a novelty claim.  In particular,
there is no assertion that $s\mapsto z+sh$ is a fine-continuous path in the
full surreal topology.  Formula \eqref{eq:Q} is the construction, and its
support certificate is what licenses it.

\begin{proposition}[Time forms under conjugacy]\label{prop:time-conjugacy}
Let $F_1,F_2$ have invertible $W_{F_1},W_{F_2}$ in $\A$.  Then
\[
 H\circ F_1=F_2\circ H
 \quad\Longleftrightarrow\quad
 H^*\Omega_{F_2}=\Omega_{F_1}
 \qquad(H\in\G_+(D)).
\]
In particular, conjugate maps have identical time-form periods on every
ordinary closed curve in $D$.
\end{proposition}
\begin{proof}
The equality of time forms reads
$H'/W_{F_2}(H)=1/W_{F_1}$, which is equivalent to the vector-field identity
in \cref{prop:intertwine}.  Period invariance now follows from
\cref{lem:homotopy}.
\end{proof}

\begin{remark}[The function class matters]
The periods here are attached to coefficient functions on $D$.  If one
allows an unrelated choice of primitive on each monad, one has enlarged the
function class and can destroy the global obstruction.  The theorem is
about a single element of $\Oh(D)((t^\Gamma))$, not a collection of
unrelated infinitesimal germs.  Likewise, a period about an ordinary
puncture is not being identified with a local residue at an arbitrary
surcomplex point.
\end{remark}

\section{Difference equations and their complete obstruction}
\label{sec:difference}

Let $F\in\G_+(D)$ be nonidentity, set $W=W_F$, and let
$\DD=W\partial_z$.  Define the scalar formal series
\begin{equation}\label{eq:E-B}
 E(X)=\frac{e^X-1}{X}=\sum_{n\geq0}\frac{X^n}{(n+1)!},\qquad
 B(X)=\frac{X}{e^X-1}=\sum_{n\geq0}\frac{B_n}{n!}X^n,
\end{equation}
where $B_0=1$ and $B_1=-1/2$.  Substitution of $\DD$ into either series
is strongly summable.  These operators are inverse automorphisms of the
underlying strongly linear space $\A$; no multiplicativity of $B(\DD)$
or $E(\DD)$ is asserted.

\begin{theorem}[Discrete equation as a primitive problem]
\label{thm:difference}
For a nonidentity $F\in\G_+(D)$,
\begin{equation}\label{eq:images}
 \ker\Delta_F=K,\qquad
 \Delta_F(\A)=W\partial_z\A.
\end{equation}
For $g\in\A$, the equation
\begin{equation}\label{eq:cohomological}
 \psi\circ F-\psi=g
\end{equation}
has a solution $\psi\in\A$ if and only if the meromorphic Hahn quotient
$g/W$ belongs to $\A$ and
\begin{equation}\label{eq:periodcondition}
 \oint_\gamma\frac{g(z)}{W(z)}\,dz=0
 \quad\text{for every ordinary closed curve }\gamma\subset D.
\end{equation}
If $P' =g/W$, all solutions are
\begin{equation}\label{eq:bernoulli-solver}
 \psi=B(\DD)P+C,\qquad C\in K.
\end{equation}
When the leading coefficient of $W$ is nowhere zero, membership $g/W\in\A$
is automatic.
\end{theorem}

\begin{proof}
By \cref{thm:logexp} and the formal identity in \eqref{eq:E-B},
\begin{equation}\label{eq:factor-difference}
 \Delta_F=e^{\DD}-I=\DD E(\DD)=E(\DD)\DD.
\end{equation}
The identities $E(\DD)B(\DD)=I=B(\DD)E(\DD)$ follow coefficientwise
from the scalar identities; \cref{cor:dependency} permits every
rearrangement.  Thus the image in \eqref{eq:images} equals $\DD\A$.
The kernel of $\DD$ is $K$, since $W\ne0$, $\A$ is a domain, and the
kernel of $\partial_z$ is $K$.  Both $E(\DD)$ and its inverse act as the
identity on constants, so \eqref{eq:factor-difference} also gives
$\ker\Delta_F=K$.

Since division by $W$ is available in $\mathcal B$, membership in
$W\partial_z\A$ is equivalent to $g/W$ being the derivative of an
element of $\A$.  \Cref{lem:primitive} gives precisely the two stated
conditions.  Finally
\[
 \Delta_F B(\DD)P=\DD P=W P'=g,
\]
and the kernel statement gives uniqueness up to $K$.  The last assertion
uses \cref{lem:unit}.
\end{proof}

\begin{remark}[Why the forcing, not a corrected forcing, supplies the periods]
One might instead multiply \eqref{eq:cohomological} by $B(\DD)$ and
attempt to integrate $B(\DD)g/W$.  This gives the same obstruction:
for $n\geq1$,
\[
 \frac{\DD^n g}{W}=(\DD^{n-1}g)'
\]
is exact.  Thus all Bernoulli correction terms have zero periods.  The
factorization \eqref{eq:factor-difference} makes the uncorrected quotient
$g/W$ and its obstruction especially transparent.
\end{remark}

\subsection{Abel functions}

\begin{corollary}[Exact global Abel criterion]\label{cor:abel}
Suppose the leading coefficient of $W_F$ is nowhere zero.  The equation
$\psi\circ F=\psi+1$ has a solution in $\A$ if and only if all the
periods of $\Omega_F$ vanish.  In this case any primitive $P'=1/W_F$
already satisfies the equation; all solutions differ from $P$ by a
constant in $K$.
\end{corollary}
\begin{proof}
Apply \cref{thm:difference} with $g=1$.  In addition,
$\DD P=1$ and $\DD^2P=0$, so
$e^{\DD}P=P+1$.  The Bernoulli formula would give
$B(\DD)P=P-1/2$, which differs only by a constant.
\end{proof}

The equality $P'=1/W_F$ guarantees a nonvanishing derivative on the halo
under the leading-unit hypothesis.  It does not make $P$ globally
injective on a multiply covered domain.  Even in ordinary complex analysis,
a locally univalent primitive need not be globally univalent.  Thus the
corollary is an Abel-function theorem, not an unconditional global
conjugacy with translation on a subset of $K$.

\subsection{An exact sequence on the punctured plane}

Let $D_r=\C\setminus\{a_1,\ldots,a_r\}$ with distinct ordinary points.
Choose positively oriented small ordinary circles $\gamma_j$ about $a_j$
which contain no other puncture.  Put
\begin{equation}\label{eq:periodmap}
 (\mathcal P_Fg)_j=\frac1{2\pi i}
       \oint_{\gamma_j}\frac{g}{W_F}\,dz.
\end{equation}

\begin{theorem}[Complete finite-dimensional discrete obstruction]
\label{thm:exact-sequence}
If $F\in\G_+(D_r)$ has a nowhere-zero leading displacement coefficient,
then the sequence \eqref{eq:exact-intro}, with the map
\eqref{eq:periodmap}, is an exact sequence of $K$-vector spaces.
A right inverse of $\mathcal P_F$ is
\begin{equation}\label{eq:rightinverse}
 (q_1,\ldots,q_r)\longmapsto
       W_F(z)\sum_{j=1}^r\frac{q_j}{z-a_j}.
\end{equation}
Consequently every forcing admits the decomposition
\begin{equation}\label{eq:forcing-decomposition}
 g=\Delta_F\psi+
       W_F\sum_{j=1}^r\frac{(\mathcal P_Fg)_j}{z-a_j},
\end{equation}
with $\psi$ determined up to a constant after the second term is fixed.
\end{theorem}

\begin{proof}
The initial kernel statement is \cref{thm:difference}.  Every ordinary
closed curve in $D_r$ has the same periods, against a holomorphic
one-form, as the integer combination of small puncture circles specified
by its winding numbers.  For completeness, one may remove small disks
about the punctures inside a large disk containing the curve, and cut the
remaining planar region along arcs joining the boundary components.  The
cut region is simply connected, and the integrals along the two sides of
each cut cancel.  This reduces its periods to those of the boundary
circles.  Applying this ordinary argument to each coefficient proves that
vanishing of the $r$ components in \eqref{eq:periodmap} is equivalent to
\eqref{eq:periodcondition}.  Thus the middle exactness follows from
\cref{thm:difference}.

The rational differential $dz/(z-a_j)$ has integral $2\pi i$ around
$\gamma_j$ and zero around the other chosen circles.  Hence
\eqref{eq:rightinverse} is a right inverse, proving surjectivity.  Subtract
this right-inverse term from $g$ and apply exactness to obtain
\eqref{eq:forcing-decomposition}.
\end{proof}

No identification of the infinite-dimensional space $\A$ with
$K\otimes_\C\Oh(D_r)$ has been used.  The finite-dimensional quotient is
computed directly by coefficientwise periods and explicit representatives.

\subsection{A local obstruction when the leading coefficient has zeros}

The quotient condition in \cref{thm:difference} must not be dropped outside
the nonsingular stratum.  On a disk containing $0$, take
\[
 W=\eps z^2,\qquad F=\frac{z}{1-\eps z},
 \qquad\eps=t^\delta,\quad\delta>0.
\]
This $F$ is the exponential of $W\partial_z$ in $\G_+(D)$.
The equation $\psi\circ F-\psi=1$ has no solution in $\A$ because
$1/W=\eps^{-1}z^{-2}$ is not holomorphic at $0$.  Equivalently, $F$ fixes
$0$, and evaluation of the equation there would give $0=1$.
On a punctured disk, by contrast, $-1/(\eps z)$ is an Abel function.
This example separates a local divisibility obstruction from the global
period obstruction.

\section{Complete global conjugacy classification}
\label{sec:classification}

Fix $\delta>0$, set $\eps=t^\delta$, and let $a\in\Oh(D)$ be nowhere
zero.  Define the stratum
\begin{equation}\label{eq:stratum}
 \G_{\delta,a}(D)=
 \{F\in\G_+(D):F=z+\eps a+r,\ \val(r)>\delta\}.
\end{equation}
By \cref{thm:logexp}, this is equivalently the class of maps whose
logarithmic vector fields have the form
$W_F=\eps(a+u)$ with $u\in\A_{>0}$.
The normalized time form is
\begin{equation}\label{eq:normalized-form}
 \beta_F=\eps\Omega_F=b_F(z)\,dz,
 \qquad b_F=a^{-1}+\text{positive Hahn terms}.
\end{equation}

The leading data $(\delta,a)$ are preserved by conjugacy in $\G_+(D)$.
Indeed, the identity $W_2(H)=H'W_1$ of
\cref{prop:intertwine} does not change the leading coefficient or valuation:
$H=z+$ positive terms and $H'=1+$ positive terms.  Thus it is natural to
classify one stratum at a time.

\subsection{A global displacement inverse}

\begin{lemma}[Implicit displacement lemma]\label{lem:displacement}
Let $b=b_0+u\in\A_{\geq0}$, where $b_0\in\Oh(D)$ is nowhere zero and
$u\in\A_{>0}$.  For each $p\in\A_{>0}$, there is a unique
$h\in\A_{>0}$ satisfying
\begin{equation}\label{eq:implicit-displacement}
 b h+\frac{b'}{2!}h^2+\frac{b''}{3!}h^3+\cdots=p.
\end{equation}
Every coefficient of $h$ is holomorphic on $D$.  If
$S=\supp(u)\cup\supp(p)$, its support is contained in
$S^*\setminus\{0\}$.  If $p(z_*)=0$, then $h(z_*)=0$.
\end{lemma}

\begin{proof}
When $S$ is empty, $h=0$ works.  Otherwise let
$M=S^*$, a well-ordered monoid by \cref{lem:neumann}.  Set the coefficient
of $h$ at $0$ to zero and determine its coefficients for
$\gamma\in M\setminus\{0\}$ in increasing order.

The coefficient of $h_\gamma$ in the left side of
\eqref{eq:implicit-displacement} is precisely $b_0$.  Any other linear
term comes from $u_\alpha h_\nu$ with $\alpha>0$ and
$\nu=\gamma-\alpha<\gamma$.  In a term containing $n\geq2$ factors
of $h$, every factor exponent is positive and strictly smaller than
$\gamma$.  The coefficient is a finite sum: the possible factors can be
expanded into words in $S$, and \cref{lem:neumann} gives finitely many
such words of total exponent $\gamma$.  Therefore the recursion has the
form
\[
 h_\gamma=b_0^{-1}(p_\gamma-R_\gamma),
\]
where $R_\gamma$ is a finite expression in earlier coefficients and
ordinary derivatives of $b$.  It defines a holomorphic function on $D$.
This constructs a Hahn series with the claimed support and proves existence.

For unrestricted uniqueness, let $h_1,h_2\in\A_{>0}$ be two solutions,
not necessarily with supports in the chosen monoid.  The difference of the
two left sides factors as
\[
 (h_1-h_2)(b_0+v),\qquad v\in\A_{>0}.
\]
This follows by factoring each difference of powers; the resulting sum is
strongly summable using the finite union of the positive supports of
$u,h_1,h_2$.  The second factor is a unit by \cref{lem:unit}.  Thus
$h_1=h_2$.

At $z_*$, evaluation gives the same scalar implicit equation over $K$,
whose linear leading coefficient is $b_0(z_*)\ne0$.  If $p(z_*)=0$,
zero is a solution and the same factorization proves it is the only
positive solution.  Hence $h(z_*)=0$.
\end{proof}

The common-domain conclusion is stronger than merely solving a separate
germ problem at each point.  Every stage divides only by the same
nowhere-zero ordinary function $b_0$ and uses derivatives of functions
already holomorphic on $D$.  No radius is lost at any Hahn exponent.

\subsection{The classification theorem}

\begin{theorem}[Complete period invariant and marked conjugacy]
\label{thm:classification}
For $F_1,F_2\in\G_{\delta,a}(D)$, the following are equivalent:
\begin{enumerate}[label=(\roman*)]
\item There is $H\in\G_+(D)$ with $H\circ F_1=F_2\circ H$.
\item For every ordinary closed piecewise $C^1$ curve $\gamma\subset D$,
\[
 \oint_\gamma\Omega_{F_1}=\oint_\gamma\Omega_{F_2}.
\]
\item $\beta_{F_1}-\beta_{F_2}=dp$ for some $p\in\A_{>0}$.
\end{enumerate}
When these conditions hold, for every specified $z_*\in D$ there is a
unique conjugacy $H$ as in \textup{(i)} satisfying $H(z_*)=z_*$.
One obtains it by taking the normalized primitive $p(z_*)=0$ in
\textup{(iii)} and solving
\begin{equation}\label{eq:construct-conjugacy}
 \sum_{n\geq1}\frac{b_{F_2}^{(n-1)}}{n!}h^n=p,
 \qquad H=z+h.
\end{equation}
\end{theorem}

\begin{proof}
By \cref{prop:time-conjugacy}, condition (i) implies (ii).
Multiplication by the nonzero constant $\eps$ shows that (ii) is equivalent
to vanishing of all periods of $\beta_{F_1}-\beta_{F_2}$.
Both normalized time forms have the same leading coefficient $a^{-1}$.
Their difference therefore has positive support.  The normalized primitive
of \cref{lem:primitive}, when it exists, also has positive support.  This
proves equivalence of (ii) and (iii), including the normalization of $p$.

Assume (iii).  The leading coefficient of $b_{F_2}$ is $a^{-1}$, which
is nowhere zero.  Apply \cref{lem:displacement} to obtain the unique
positive $h$ in \eqref{eq:construct-conjugacy}.  The explicit pullback
identity of \cref{lem:homotopy} gives
\[
 H^*\beta_{F_2}-\beta_{F_2}=dp
             =\beta_{F_1}-\beta_{F_2}.
\]
Thus $H^*\beta_{F_2}=\beta_{F_1}$, equivalently
$H^*\Omega_{F_2}=\Omega_{F_1}$.  By
\cref{prop:time-conjugacy}, $H$ is a conjugacy.  Its inverse belongs to
$\G_+(D)$ by \cref{thm:logexp}, and
\cref{lem:displacement} gives $H(z_*)=z_*$.

Finally suppose $\widetilde H=z+\widetilde h$ is any conjugacy fixing
$z_*$.  The same pullback identity yields
\[
 d\bigl(Q_{\widetilde h}(b_{F_2})-p\bigr)=0.
\]
The expression in parentheses is a constant in $K$, and its value at
$z_*$ is zero.  Hence $\widetilde h$ satisfies the same implicit equation
\eqref{eq:construct-conjugacy}.  Its uniqueness in
\cref{lem:displacement} proves $\widetilde H=H$.
\end{proof}

\begin{corollary}[Simply connected domains]\label{cor:simplyconnected}
If $D$ is simply connected, every two maps in $\G_{\delta,a}(D)$ are
conjugate by a unique near-identity change of variable fixing a prescribed
ordinary basepoint.  In particular, each is conjugate to
$\exp(\eps a\partial_z)z$.
\end{corollary}
\begin{proof}
Every ordinary holomorphic coefficient one-form on $D$ has a primitive,
so condition (iii) of \cref{thm:classification} holds.
\end{proof}

For $a=1$, the reference map is $z+\eps$ and this is a genuine global
near-identity conjugacy to that translation on $D^\#$.  For general $a$,
the reference map is the indicated vector-field exponential; no
univalence of a primitive of $dz/a$ is required.

\subsection{Why the hypotheses cannot be silently weakened}

Three restrictions do actual work.  The leading coefficient $a$ must not
vanish, because inversion and the implicit displacement recursion divide by
$a$ or $a^{-1}$.  The conjugating map must have identity ordinary part,
because otherwise it can act nontrivially on ordinary homology and on the
leading form.  Finally, all coefficient functions must be defined on the
same $D$; replacing global functions by independent germs removes the
meaning of a single period comparison.  The theorem is not claimed after
any of these hypotheses is erased.

\section{Explicit moduli and the centralizer}
\label{sec:moduli}

\subsection{Every period vector occurs}

Continue with $D_r=\C\setminus\{a_1,\ldots,a_r\}$.  For the fixed
stratum $(\delta,a)$ define
\begin{equation}\label{eq:moduli-map}
 q_j(F)=\frac1{2\pi i}\oint_{\gamma_j}
          \left(\beta_F-\frac{dz}{a(z)}\right).
\end{equation}
Its values belong to $\mm_K$, because the difference has positive support.

\begin{theorem}[An explicit moduli space]\label{thm:moduli}
The map $F\mapsto(q_1(F),\ldots,q_r(F))$ induces a bijection
\begin{equation}\label{eq:moduli-bijection}
 \G_{\delta,a}(D_r)/\text{conjugacy by }\G_+(D_r)
       \ \longleftrightarrow\ \mm_K^r.
\end{equation}
For $\mathbf q=(q_1,\ldots,q_r)\in\mm_K^r$, a representative is
\begin{equation}\label{eq:moduli-representative}
 b_{\mathbf q}(z)=\frac1{a(z)}+\sum_{j=1}^r\frac{q_j}{z-a_j},
 \qquad W_{\mathbf q}=\frac{\eps}{b_{\mathbf q}},
 \qquad F_{\mathbf q}=\exp(W_{\mathbf q}\partial_z)z.
\end{equation}
\end{theorem}

\begin{proof}
The finite union of the supports of the $q_j$ is positive and well ordered.
Thus $b_{\mathbf q}\in\A_{\geq0}$ with nowhere-zero leading coefficient
$a^{-1}$, and it is invertible by \cref{lem:unit}.  Consequently
$W_{\mathbf q}=\eps(a+\text{positive terms})$, and
$F_{\mathbf q}$ belongs to the specified stratum.  Its normalized time
form is exactly $b_{\mathbf q}dz$.  The integrals of the simple-pole
representatives give $q_j(F_{\mathbf q})=q_j$.

Conjugate maps have the same vector by period invariance.  Conversely, equal
vectors give equal periods around each puncture.  As in
\cref{thm:exact-sequence}, this gives equality of every ordinary period.
\Cref{thm:classification} then gives a conjugacy.  This proves both
surjectivity and injectivity in \eqref{eq:moduli-bijection}.
\end{proof}

The parametrization is algebraic over the Hahn constants in its displayed
one-form representatives.  It imposes no countability, finite generation,
or rank-one condition on the possible supports of the $q_j$.

\subsection{Centralizers and uniqueness of roots}

The next result needs no nowhere-zero leading coefficient.

\begin{theorem}[Centralizer of a nonidentity infinitesimal map]
\label{thm:centralizer}
Let $F\in\G_+(D)$ be nonidentity, with $W=W_F$ and
$\delta=\val(W)>0$.  Its centralizer in $\G_+(D)$ is
\begin{equation}\label{eq:centralizer}
 Z_{\G_+(D)}(F)=
 \left\{\exp(cW\partial_z)z:
 c\in K,\ c=0\text{ or }\val(c)+\delta>0\right\}.
\end{equation}
The parametrization is an isomorphism from the indicated additive subgroup
of $K$ to the centralizer.  Every $F\in\G_+(D)$ has exactly one $n$th
root in $\G_+(D)$ for every integer $n\geq1$, namely
$\exp((W_F/n)\partial_z)z$.  The group $\G_+(D)$ is torsion-free.
\end{theorem}

\begin{proof}
If $H$ commutes with $F$, their substitution operators commute, so their
summable logarithms commute.  Write $W_H=V$.  The bracket is
\[
 [W\partial_z,V\partial_z]=(WV'-VW')\partial_z,
\]
and therefore $WV'-VW'=0$.  In the meromorphic Hahn field $\mathcal B$,
\[
 (V/W)'=0.
\]
Its constant field is $K$, so $V=cW$ for some $c\in K$.
If $c\ne0$, positivity of $V$ is exactly
$\val(c)+\delta>0$.  Conversely, for each such $c$, the positive vector
field $cW$ exponentiates, and its exponential commutes with that of $W$.
The group law follows from
\[
 \exp(cW\partial_z)\exp(dW\partial_z)
     =\exp((c+d)W\partial_z).
\]
Injectivity follows by taking logarithms and using $W\ne0$.

The displayed $n$th root certainly exists.  If $H^n=F$ and $W_H=V$,
then $T_H^n=\exp(nV\partial_z)$; uniqueness of logarithms gives
$nV=W_F$.  This proves uniqueness.  If $H^n=\id$, the same argument gives
$nW_H=0$, hence $H=\id$ in characteristic zero.
\end{proof}

If the leading coefficient of $W_F$ is nowhere zero at $z_*$, no
nonidentity centralizer element fixes $z_*$.  Indeed, its leading
displacement is the nonzero value of $cW_F$ there.  Thus the basepoint
normalization in \cref{thm:classification} removes exactly the
centralizer freedom.  More generally, the unnormalized conjugacies between
two conjugate maps form a torsor under either centralizer.

The time constants in \eqref{eq:centralizer} may have negative valuation,
provided $cW$ remains positive.  This is a legitimate Hahn-defined family
of maps, not a claim that arbitrary large times are available for every
flow.  The positivity threshold is part of the conclusion.

\section{Nonconjugate dynamics beyond every one-scale jet}
\label{sec:flat}

\subsection{The flat ideal of a noncofinal scale}

Fix $\delta>0$, set $\eps=t^\delta$, and define
\begin{equation}\label{eq:flat-ideal}
 J_\delta=\{0\}\cup
 \{c\in K:\val(c)>n\delta\text{ for every integer }n\geq1\}.
\end{equation}
It is an ideal in the valuation ring $K_{\geq0}$.  It is nonzero precisely
when there is some $\eta\in\Gamma$ larger than all the $n\delta$.
The forward implication follows by taking the valuation of a nonzero
member; the reverse implication follows from $t^\eta\in J_\delta$.
In particular $J_\delta=0$ when the multiples of $\delta$ are cofinal,
as happens in an Archimedean value group.

For coefficient functions, the corresponding ideal is
\begin{equation}\label{eq:flat-functions}
 I_\delta(D)=\bigcap_{n\geq1}\eps^n\A_{\geq0}
   =\{0\}\cup\{f\in\A:\val(f)>n\delta\text{ for all }n\geq1\}.
\end{equation}
The equivalence between weak and strict inequalities here follows by
replacing $n$ with $n+1$.  It is essential that these ideals are formed in
$\A_{\geq0}$, \emph{not} in $\A$: $\eps$ is a unit in the latter,
so its powers generate the whole ring there.

Consequently the map to one-scale jets
\begin{equation}\label{eq:jet-map}
 \A_{\geq0}\longrightarrow
       \varprojlim_n\A_{\geq0}/\eps^n\A_{\geq0}
\end{equation}
has kernel $I_\delta(D)$.  In the noncofinal case it loses nonzero
information.  The next theorem shows that it can lose a complete family
of global conjugacy obstructions, not merely some irrelevant constants.

\subsection{An exact family of invisible moduli}

\begin{theorem}[Flatly indistinguishable but globally nonconjugate]
\label{thm:flat}
Let $r\geq1$, let $D_r=\C\setminus\{a_1,\ldots,a_r\}$, and suppose
$J_\delta\ne0$.  For $\mathbf c=(c_1,\ldots,c_r)\in J_\delta^r$ set
\begin{equation}\label{eq:flat-family}
 R_{\mathbf c}(z)=\sum_{j=1}^r\frac{c_j}{z-a_j},\qquad
 W_{\mathbf c}=\frac{\eps}{1+\eps R_{\mathbf c}},\qquad
 F_{\mathbf c}=\exp(W_{\mathbf c}\partial_z)z.
\end{equation}
Then all these maps belong to $\G_{\delta,1}(D_r)$ and satisfy:
\begin{enumerate}[label=(\alph*)]
\item $F_{\mathbf0}=z+\eps$ and
$F_{\mathbf c}-F_{\mathbf0}\in I_\delta(D_r)$ for every $\mathbf c$.
If $\mathbf c\ne0$ and $\eta=\min_{c_j\ne0}\val(c_j)$, then more precisely
\begin{equation}\label{eq:flat-valuation}
 \val(F_{\mathbf c}-F_{\mathbf0})=2\delta+\eta.
\end{equation}
\item Their time forms are exactly
\begin{equation}\label{eq:exact-flat-form}
 \Omega_{F_{\mathbf c}}=\eps^{-1}dz+
                    \sum_{j=1}^r\frac{c_j\,dz}{z-a_j}.
\end{equation}
\item $F_{\mathbf c}$ and $F_{\mathbf d}$ are conjugate by an element
of $\G_+(D_r)$ if and only if $\mathbf c=\mathbf d$.
\item A global Abel function in $\A_\Gamma(D_r)$ exists for
$F_{\mathbf c}$ if and only if $\mathbf c=\mathbf0$.
\end{enumerate}
\end{theorem}

\begin{proof}
The finite union of the supports of the $c_j$ is well ordered and positive.
Hence $R_{\mathbf c}\in\A_{>0}$.  The denominator in
\eqref{eq:flat-family} has leading coefficient $1$ and is a unit, so
$W_{\mathbf c}=\eps+$ terms of valuation greater than $\delta$.
The exponential theorem places $F_{\mathbf c}$ in the stated stratum.
For $\mathbf c=0$ the vector field is constant, and its exponential is
exactly $z+\eps$.

If $\mathbf c\ne0$, then
$\val(R_{\mathbf c})=\eta$.  Indeed, a nontrivial linear combination of
the rational functions $(z-a_j)^{-1}$ cannot vanish identically: taking
ordinary residues at the punctures recovers its coefficients.
The geometric expansion gives
\begin{equation}\label{eq:Wflat-difference}
 W_{\mathbf c}-\eps=-\eps^2R_{\mathbf c}
                    +\eps^3R_{\mathbf c}^2-\cdots,
 \qquad\val(W_{\mathbf c}-\eps)=2\delta+\eta.
\end{equation}
Put $V=W_{\mathbf c}-\eps$ and $\lambda=2\delta+\eta$.
In the exponential expansion of $(\eps\partial_z+V\partial_z)^n z$
for $n\geq2$, the word with no factor $V\partial_z$ vanishes, because
$(\eps\partial_z)^2z=0$.  Every other word has valuation at least
$\lambda+(n-1)\delta$, since derivatives do not lower Hahn valuation.
This is strictly greater than $\lambda$.  All such words are summable
by \cref{lem:neumann}.  Thus the leading difference between
$F_{\mathbf c}$ and $z+\eps$ comes from $V$ itself, proving
\eqref{eq:flat-valuation}.  Since $\eta>n\delta$ for all $n$, the
difference belongs to \eqref{eq:flat-functions}.

The logarithm of $F_{\mathbf c}$ is $W_{\mathbf c}$, by the inverse
identity in \cref{thm:logexp}.  Taking its reciprocal gives the exact
formula \eqref{eq:exact-flat-form}.  Thus $(2\pi i)^{-1}\oint_{\gamma_j}\Omega_{F_{\mathbf c}}=c_j$;
for the normalized time form $\beta_{F_{\mathbf c}}$ the corresponding
value is $\eps c_j$.
Thus equality of time-form periods is equivalent to
$\mathbf c=\mathbf d$.  \Cref{thm:classification} proves (c), while
\cref{cor:abel} proves (d).
\end{proof}

\begin{remark}[A failure of separated one-scale information, not of Hahn equality]
All the maps in the theorem have the same image under the one-scale jet map
\eqref{eq:jet-map}.  They do \emph{not} have the same full Hahn coefficients.
The missing coefficients occur at exponents larger than every $n\delta$.
Thus the result does not contradict uniqueness of Hahn expansions, nor
separation by arbitrary value-group cutoffs.  It shows exactly why replacing
arbitrary-rank support control by a single ``small parameter completion''
can erase a global obstruction.
\end{remark}

\subsection{A concrete surcomplex instance}

Take the ordered group
\[
 \Gamma=\Q\oplus\Q\quad\text{with lexicographic order},\qquad
 \delta=(0,1),\quad\eta=(1,0).
\]
For every integer $n\geq1$, $\eta>n\delta$, because comparison is decided
by the first coordinate.  This group can be represented inside the
surreals as $\Q\omega+\Q$, with $(q,r)\mapsto q\omega+r$.
Under the normal-form embedding,
\[
 \eps=\omega^{-1},\qquad c=t^\eta=\omega^{-\omega}.
\]
On $D=\C^*$, the field
\begin{equation}\label{eq:surreal-flat-example}
 W_c(z)=\frac{\omega^{-1}}{1+\omega^{-(\omega+1)}/z}
\end{equation}
therefore defines an explicit surcomplex map $F_c=\exp(W_c\partial_z)z$.
It has the same truncation as $z+\omega^{-1}$ modulo every positive integer
power of $\omega^{-1}$, but
\[
 \frac1{2\pi i}\oint_{|z|=1}\Omega_{F_c}=\omega^{-\omega}\ne0.
\]
It is not globally near-identity conjugate to that translation and has no
single-valued Hahn-coherent Abel function on the punctured plane.

This is a global effect at a strictly smaller scale, not a claim about
ordinary exponentially small asymptotics of complex numbers.  The quantity
$\omega^{-\omega}$ is an exact Hahn monomial in a rank-two workspace.
A nonzero ordinary multiple of it already gives a different member of the
family; the full parameter space $J_\delta$ is much larger.

\section{Worked dynamical examples}
\label{sec:examples}

\subsection{First terms of the logarithm}

For a map written in powers of one positive monomial,
\[
 F=z+\eps v+\eps^2b+\eps^3c+\Oe4,
\]
where $v,b,c\in\Oh(D)$, direct expansion of
\eqref{eq:operatorlog} gives
\begin{align}
 W_F={}&\eps v+\eps^2\left(b-\tfrac12vv'\right)\notag\\
 &+\eps^3\left(c-\tfrac12(vb'+bv')
       +\tfrac1{12}v^2v''+\tfrac13v(v')^2\right)+\Oe4.
 \label{eq:log-first-terms}
\end{align}
Here $\Oe4$ means membership in $\eps^4\Oh(D)[[\eps]]$; this notation
is used only in examples already specified to be one-parameter series.
It is not used as a substitute for arbitrary-rank support conditions.

To obtain the cubic terms, write $h=F-z$.  Up to cubic perturbative degree,
\[
 \Delta z=h,\quad
 \Delta^2z=hh'+\tfrac12h^2h''+\cdots,\quad
 \Delta^3z=h(h')^2+h^2h''+\cdots.
\]
Inserting the logarithm coefficients $1,-1/2,1/3$ gives
$h-\frac12hh'+\frac1{12}h^2h''+\frac13h(h')^2$, from which
\eqref{eq:log-first-terms} follows.

If $v$ is nowhere zero, its time form begins
\begin{equation}\label{eq:time-first-terms}
 \Omega_F=
 \left[\frac1{\eps v}+\frac{v'}{2v}-\frac{b}{v^2}
                 +O(\eps)\right]dz.
\end{equation}
Thus already at the next order, the ordinary logarithmic derivative of the
leading vector field can create a nonzero period.

\subsection{Euler maps and an exact iterative-residue obstruction}

Consider, on $D=\C^*$,
\begin{equation}\label{eq:euler-map}
 F_p(z)=z+\eps z^{p+1},\qquad p\geq1.
\end{equation}
The support is contained in $\N\delta$, so these examples are available
in every group containing $\delta>0$.

\begin{proposition}[Exact residue of the Euler time form]
\label{prop:euler-residue}
The logarithmic vector field and its reciprocal have the forms
\begin{align}
 W_{F_p}(z)&=\sum_{n\geq1}c_{p,n}\eps^n z^{np+1},
      &c_{p,1}&=1,& c_{p,2}&=-\frac{p+1}{2},\label{eq:euler-grading}\\
 \frac1{W_{F_p}(z)}&=\eps^{-1}z^{-p-1}
       +\frac{p+1}{2}z^{-1}
       +\sum_{n\geq1}d_{p,n}\eps^n z^{np-1}.
       \label{eq:euler-reciprocal}
\end{align}
In particular,
\begin{equation}\label{eq:euler-residue}
 \Res_{z=0}\Omega_{F_p}=\frac{p+1}{2}
\end{equation}
exactly, with no omitted positive-valuation correction.  There is no global
Abel function for $F_p$ in $\Oh(\C^*)((t^\Gamma))$.
\end{proposition}

\begin{proof}
Taylor substitution by \eqref{eq:euler-map} preserves the grading in which
$\eps^n z^{np+1}$ has grade $n$.  Explicitly,
\[
 (z+\eps z^{p+1})^{kp+1}
     =z^{kp+1}(1+\eps z^p)^{kp+1}.
\]
Every term of the finite difference and its iterates therefore has the
form prescribed in \eqref{eq:euler-grading}.  Taking the logarithm preserves
it.  The first two coefficients follow from
\eqref{eq:log-first-terms}.

Factor $W_{F_p}=\eps z^{p+1}C_p(\eps z^p)$ with
$C_p(u)=1-(p+1)u/2+\cdots$.  Inverting $C_p$ gives
\eqref{eq:euler-reciprocal}.  For $n\geq1$, the exponent $np-1$ is
nonnegative.  Therefore none of the last terms has a $z^{-1}$ coefficient,
and \eqref{eq:euler-residue} is exact.  The period about $0$ is nonzero,
so \cref{cor:abel} forbids a global Abel function.
\end{proof}

The leading vector field $\eps z^{p+1}$ itself has zero period on
$\C^*$, and its exponential has Abel function
$-1/(p\eps z^p)$.  Hence passing from that exact flow to its Euler map
changes the global obstruction even though their first displacements agree.
The constant in \eqref{eq:euler-residue} is the familiar iterative-residue
phenomenon; the assertion here is its exact interpretation in the stated
common-domain Hahn category, not the discovery of that classical invariant.

For $p=1$ the first terms are
\begin{equation}\label{eq:euler-known}
 W_{F_1}=\eps z^2-\eps^2z^3+\tfrac32\eps^3z^4
 -\tfrac83\eps^4z^5+\tfrac{31}{6}\eps^5z^6
 -\tfrac{157}{15}\eps^6z^7+\Oe7.
\end{equation}
These coefficients agree with the standard modified equation in Hairer's
example \cite{hairer}, after multiplication by the step size.  They are
recomputed in the accompanying verifier.

\begin{corollary}[A global coherent normal form for the Euler family]
\label{cor:euler-normal}
The map $F_p$ is conjugate on $(\C^*)^\#$, by a near-identity
Hahn-coherent change of variable, to the time-one map of
\begin{equation}\label{eq:euler-normal}
 V_p(z)=\frac{\eps z^{p+1}}
             {1+\frac{p+1}{2}\eps z^p}.
\end{equation}
The conjugacy fixing $z=1$ is unique.  It is not a conjugacy of $F_p$ to
the time-one map of $\eps z^{p+1}$.
\end{corollary}
\begin{proof}
The time form of \eqref{eq:euler-normal} is exactly
\[
 \eps^{-1}z^{-p-1}dz+\frac{p+1}{2}\frac{dz}{z}.
\]
It has the same period as \eqref{eq:euler-reciprocal}.
The difference consists of nonnegative powers of $z$, coefficientwise,
so it is exact on $\C^*$.  Apply \cref{thm:classification}.
The last assertion follows because the leading vector-field flow has zero
period while $F_p$ does not.
\end{proof}

On any simply connected slit domain in $\C^*$, the normal form has the
local Abel function
\[
 -\frac1{p\eps z^p}+\frac{p+1}{2}\log z.
\]
Its logarithmic monodromy is precisely the obstruction to a single-valued
global Abel function on $\C^*$.

\subsection{A simple difference equation with an explicit correction}

Take $W=\eps z$, so $F=e^\eps z$, on $D=\C^*$.  The forcing $g=z^m$
with nonzero integer $m$ has zero period after division by $W$, and
\[
 \psi(z)=\frac{z^m}{e^{m\eps}-1}
\]
is a solution.  The denominator has valuation $\delta$, so its reciprocal
is a legitimate Laurent-Hahn constant.  Expanding it gives
\[
 \psi(z)=\left(\frac1{m\eps}-\frac12
                    +\frac{m\eps}{12}-\cdots\right)z^m,
\]
exactly the Bernoulli solver of \cref{thm:difference} applied to
$P=z^m/(m\eps)$.  For the constant forcing $g=1$, on the other hand,
$g\,dz/W=\eps^{-1}dz/z$ has nonzero period and no global solution exists.
This directly displays the one-dimensional cokernel for the punctured plane.

\subsection{Local straightening does not erase global periods}

The maps $F(z)=(1+\eps)z$ and $G(z)=e^\eps z$ have leading displacement
$\eps z$ on $\C^*$.  Their logarithmic vector fields are respectively
$\log(1+\eps)z$ and $\eps z$.  Their periods are
\[
 \frac{2\pi i}{\log(1+\eps)}\quad\text{and}\quad\frac{2\pi i}{\eps},
\]
which differ, so there is no global near-identity conjugacy on $\C^*$.
On a simply connected domain with a branch of $\log z$, there is one:
\[
 H(z)=z\exp\left(\left[\frac{\eps}{\log(1+\eps)}-1\right]\log z\right).
\]
The bracketed coefficient has positive valuation, so the exponential is a
strongly summable Hahn expansion with coefficients holomorphic on the chosen
branch domain.  Direct substitution gives $H\circ F=G\circ H$.
The failure to make that formula single-valued on $\C^*$ is exactly what
the period theorem detects.

\section{Ordinary slow time on one fixed holomorphic flow chart}
\label{sec:slowtime}

The suspension in \cref{thm:logexp} concerns an infinitesimal vector field.
One may also ask what happens on the larger time scale at which its leading
ordinary vector field moves a finite distance.  A direct exponential is
generally not a strongly summable Hahn expression on that scale.  The
correct construction first solves the leading ordinary flow and then lifts
all positive Hahn coefficients along it.

Write
\begin{equation}\label{eq:slow-vector}
 W=\eps X,\qquad
 X(z)=a(z)+\sum_{\beta\in S}t^\beta X_\beta(z),
\end{equation}
where $S\subset\Gamma_{>0}$ is well ordered and all coefficient functions
are holomorphic on $D$.  In this section $a$ may have zeros.  Let
$\Sigma\subset\C$ be an ordinary disk centered at $0$, and let
$D_0\subset D$ be an ordinary open domain.  Assume that a holomorphic
ordinary flow chart
\begin{equation}\label{eq:ordinary-flow}
 \phi:\Sigma\times D_0\longrightarrow D,
 \qquad \partial_s\phi(s,z)=a(\phi(s,z)),\qquad\phi(0,z)=z,
\end{equation}
is specified.  Set $J(s,z)=\partial_z\phi(s,z)$ and assume $J$ never
vanishes on this chart.  The latter also follows from the usual variational
equation on a nonsingular ordinary flow chart; we include it among the
explicit hypotheses used in the formula below.

\begin{theorem}[Common-chart Hahn lifting of the slow flow]
\label{thm:slowtime}
There is exactly one series
\begin{equation}\label{eq:slow-solution}
 \Psi(s,z)=\phi(s,z)+
    \sum_{\gamma\in S^*\setminus\{0\}}t^\gamma\psi_\gamma(s,z)
\end{equation}
with coefficients holomorphic on the \emph{same} product domain
$\Sigma\times D_0$, satisfying
\begin{equation}\label{eq:slow-equation}
 \partial_s\Psi=X(\Psi),\qquad\Psi(0,z)=z.
\end{equation}
Here $X(\Psi)$ means Taylor expansion at the ordinary point $\phi(s,z)$.
Uniqueness also holds among all positive-Hahn perturbations of $\phi$ on
this product domain, not just those with the displayed support bound.

Each coefficient is obtained by one ordinary integral:
\begin{equation}\label{eq:variation-formula}
 \psi_\gamma(s,z)=J(s,z)
       \int_0^s J(u,z)^{-1}R_\gamma(u,z)\,du,
\end{equation}
where $R_\gamma$ is a finite expression in previously determined
coefficients and in ordinary derivatives of the $X_\beta$ and $a$
along $\phi$.

Finally, Taylor evaluation at $s=\eps$ is strongly summable and satisfies
\begin{equation}\label{eq:slow-epsilon}
 \Psi(\eps,z)=\exp(W\partial_z)z\quad(z\in D_0^\#).
\end{equation}
More generally, for $c\in K$ with $c=0$ or $\val(c)+\delta>0$,
Taylor evaluation at $s=c\eps$ is defined and gives
$\Psi(c\eps,z)=\exp(cW\partial_z)z$.
\end{theorem}

\begin{proof}
Let $M=S^*$ and seek $q=\Psi-\phi$ with support in
$M\setminus\{0\}$.  Expand
\begin{align*}
 X(\phi+q)
 &=a(\phi)+a'(\phi)q+
      \sum_{n\geq2}\frac{a^{(n)}(\phi)}{n!}q^n\\
 &\quad+\sum_{\beta\in S}t^\beta
            \sum_{n\geq0}\frac{X_\beta^{(n)}(\phi)}{n!}q^n.
\end{align*}
All the compositions with $\phi$ are holomorphic on the specified product
because its image lies in $D$.  For each $\gamma\in M\setminus\{0\}$,
the coefficient equation is
\begin{equation}\label{eq:variational}
 \partial_s\psi_\gamma=a'(\phi)\psi_\gamma+R_\gamma,
 \qquad\psi_\gamma(0,z)=0.
\end{equation}
Every factor in $R_\gamma$ involves a positive exponent smaller than
$\gamma$, except for specified coefficients of $X$.  The number of
contributions is finite by the word lemma, exactly as in
\cref{lem:displacement}.  Thus well-ordered recursion is possible.

Differentiating the ordinary flow equation in $z$ gives
\[
 \partial_s J=a'(\phi)J,\qquad J(0,z)=1.
\]
Variation of constants solves \eqref{eq:variational} by
\eqref{eq:variation-formula}.  Since $J$ is nowhere zero and $\Sigma$ is
simply connected, the integral is independent of its path in $\Sigma$.
It is holomorphic jointly in $(s,z)$; this follows locally by integrating
the holomorphic power series in $s$, with coefficients holomorphic in $z$,
and patching the uniquely normalized primitives.  Therefore every newly
constructed coefficient is holomorphic on the original product domain.
There is no successive domain shrinkage and no requirement of a uniform
bound on the size of these coefficients as $\gamma$ varies.

The support of the resulting series is contained in the well-ordered set
$M$.  The expansion of $X(\Psi)$ is strongly summable by
\cref{lem:neumann}, so the coefficient construction indeed solves
\eqref{eq:slow-equation}.  For uniqueness, compare any two positive-Hahn
perturbations of $\phi$ satisfying that equation.  At the least exponent
where they differ, their difference solves the homogeneous version of
\eqref{eq:variational} with zero initial condition.  It is zero, a
contradiction.  This argument does not require either support to lie in
$M$.

For the last assertion, expand every coefficient at $s=0$ and substitute
$s=\eps$.  The possible exponents lie in the monoid generated by
$S\cup\{\delta\}$, and a fixed exponent has finitely many word
representations.  Hence this is a strongly summable evaluation.  Repeated
differentiation of \eqref{eq:slow-equation}, or induction using the chain
rule, gives the coefficientwise Taylor identities
\[
 \partial_s^n\Psi(0,z)=(X\partial_z)^n z\quad(n\geq0).
\]
Therefore
\[
 \Psi(\eps,z)=\sum_{n\geq0}\frac{\eps^n}{n!}
                       (X\partial_z)^n z
             =\exp(W\partial_z)z.
\]
For $s=c\eps\in\mm_K$, replace the singleton support $\{\delta\}$
by the well-ordered positive support of $c\eps$.  The same word argument
justifies evaluation, and the identical Taylor calculation gives the
stated expression with $cW$.
\end{proof}

\subsection{What this construction adds}

At an ordinary time $s_0\in\Sigma$, evaluation of
\eqref{eq:slow-solution} produces an element with leading ordinary part
$\phi(s_0,z)$ and positive Hahn corrections.  The ordinary coefficient
integrals have already been performed before this evaluation.  This is a
separate analytic operation from an unregrouped Hahn exponential.

For example, take $X=z^2$ and $W=\eps z^2$.  The slow flow is the ordinary
rational function
\[
 \phi(s,z)=\frac{z}{1-sz}
\]
on any product chart avoiding $sz=1$.  At a nonzero ordinary $s$, the raw
formal terms of
\[
 \exp(s z^2\partial_z)z=\sum_{n\geq0}s^n z^{n+1}
\]
all have Hahn exponent $0$, and infinitely many are nonzero.  They are
\emph{not} a strongly summable Hahn family.  The rational function is
available from ordinary complex analysis on its chart, not from silently
declaring that family to be a Hahn sum.

Thus \cref{thm:slowtime} licenses a particular passage to the time scale
$\eps^{-1}$: solve the leading ordinary motion and then integrate each
positive correction along it.  It does not license arbitrary exchange of
ordinary summation with Hahn summation, global continuation across a
singularity of $\phi$, or an interpretation of every surreal-valued time
as an actual continuous trajectory.

\section{Scope, proof dependencies, and further directions}
\label{sec:scope}

\subsection{The precise category in which the results hold}

The input to every main theorem consists of set-sized supports, ordinary
holomorphic coefficient functions on a named common domain, and positive
Hahn perturbations.  All differentiation is with respect to the ordinary
coordinate $z$ extended coefficientwise.  Every constant from $K$ is
annihilated.  All loops are ordinary complex loops and all their integrals
are coefficientwise.  No fine-topological convergence theorem on
$\No[i]$ is used.

In particular, these arguments do not identify
$\Oh(D)((t^\Gamma))$ with a ring of coefficient germs without a common
domain, or with $\C[[z]]((t^\Gamma))$.  They do not construct a global
sheaf from a uniform-support convention on arbitrary covers.  None of
these extra identifications is needed: the proofs stay on a specified $D$
and use support-preserving primitives there.

The entire argument is compatible with adjoining further Hahn exponents:
the displayed formulas remain valid in a larger ordered workspace, and an
equality or a nonzero existing period remains an equality or a nonzero
period under the natural embedding.  The statements still quantify over
the explicitly specified coefficient-function category in that workspace,
not over all conceivable surcomplex functions.

\subsection{Where the new work sits in the dependency chain}

The proof dependencies are
\[
 \begin{gathered}
 \text{finite-word support lemma}
 \ \Longrightarrow\ \text{Taylor calculus and positive operators}
 \\
 \Longrightarrow\ \text{canonical vector field}\
 \Longrightarrow\ \text{time form and discrete equation factorization},
 \\
 \text{primitive criterion}+
 \text{implicit displacement lemma}
 \ \Longrightarrow\ \text{complete global conjugacy},
 \\
 \text{conjugacy theorem}+
 \text{explicit rational time forms}
 \ \Longrightarrow\ \text{moduli and flat nonconjugacy}.
 \end{gathered}
\]
The slow-time theorem branches directly from the support lemma and the
ordinary variational equation.  It does not assume the classification.
The positive-operator logarithm is established theory in a new concrete
setting; the global obstruction and moduli statements are the principal
proposed contributions.

\subsection{Problems not settled by these results}

When the leading coefficient has zeros, the discrete equation remains
controlled by \cref{thm:difference}, but the conjugacy classification needs
more information.  Local zero divisors, multiplicities, and singular
normal-form obstructions must be combined with global periods.  The unit
hypothesis in the displacement lemma is exactly the point at which our
classification proof stops applying.

A conjugacy with nonidentity ordinary part can change the leading vector
field and the homology basis.  Extending the classification to that group
would require the ordinary automorphism action on the leading form and on
its periods; equal period vectors in a fixed basis would no longer be the
correct equivalence relation by themselves.

In several variables, there is no single reciprocal time one-form which
encodes a vector field, and commuting vector fields need not be constant
multiples of one another.  The one-variable centralizer proof and
finite-period normal form should therefore not be transported unchanged.

Finally, ordinary analytic specialization of a Hahn or formal parameter
requires growth estimates or an additional summation theory.  The present
results neither supply those estimates nor remove sectorial analytic
moduli.  Establishing which members of the global Hahn normal forms admit
compatible ordinary analytic realizations is a distinct problem.

\section{Reproducibility and verification record}
\label{sec:verification}

The accompanying \texttt{verify.py} works in the finite jet ring
$\Q(z)[\eps]/(\eps^8)$, except for explicitly marked exact rational
identities and finite ordered-group illustrations.  It implements Taylor
substitution, powers of the finite-difference operator, the logarithm,
vector-field exponentiation, and the Bernoulli operator using exact
arithmetic.

The recorded run, with Python 3.13.5 and SymPy 1.14.0, passes
\textbf{245 exact assertions}.  It checks the Euler families $p=1,2,3$,
logarithm--exponential round trips, inverse maps, Julia's identity, the
homogeneous grading and residue coefficients, a mixed-coefficient example,
Leibniz's identity, the Bernoulli difference solver, flow composition, a
square-root formula, a nontrivial conjugacy and time-form pullback, and
the rational identity for the flat family.  Fifty assertions are explicitly
labelled finite illustrations of the rank-two comparison; the all-integer
inequality is proved by lexicographic order in \cref{sec:flat}, not by
those tests.

The verifier handles a loss of one known coefficient on division by
$\eps$ explicitly: any subsequent reciprocal/pullback comparison stops
before the uncomputed top coefficient.  It uses no floating point tolerance.
The output is retained in \texttt{verification.txt}.

These tests are not a proof of arbitrary-rank summability, global primitive
existence, the classification, or the slow-time theorem.  Those depend on
the arguments in the main text.  A passing finite test is evidence against
an algebraic transcription error, not a replacement for a proof and not a
formal proof-assistant certificate.

\appendix
\section{A coefficient-level recipe for constructing a conjugacy}
\label{app:recipe}

The construction in \cref{thm:classification} is explicit even when it is
not an effective algorithm for arbitrary input coefficients.  Here is its
coefficient-level content.  Starting with $F_1,F_2$ in the same nonsingular
stratum, first form their logarithmic vector fields by
\eqref{eq:operatorlog}; compute the normalized reciprocal forms
$\beta_i=b_i dz$ by the geometric unit expansion; and test the periods of
$(b_1-b_2)dz$.  On a finitely punctured plane this is a finite list of Hahn
periods, though computing an arbitrary coefficient integral is not assumed
to be decidable.

When the periods vanish, integrate coefficientwise from $z_*$ to get
$p$ with $p'=b_1-b_2$ and $p(z_*)=0$.  Let $b_2=b_0+u$ and let $S$
be the positive union of the supports of $p$ and $u$.  The desired
$h=H-z$ has support in $S^*\setminus\{0\}$.  For each $\gamma$ in that
well-ordered set, define
\begin{equation}\label{eq:explicit-recursion}
 h_\gamma=b_0^{-1}\left(
 p_\gamma-[t^\gamma]\left[
 uh+\sum_{n\geq2}\frac{b_2^{(n-1)}}{n!}h^n
 \right]\right).
\end{equation}
The expression in brackets uses only earlier coefficients of $h$.
The number of operations at each exponent is finite, but the total support
may be transfinite and need not have a finite presentation.

For an ordinary single-parameter expansion
$b_2=b_0+\eps b_1+\eps^2b_2^{[2]}+\cdots$ and
$p=\eps p_1+\eps^2p_2+\cdots$, the first coefficients are
\begin{align*}
 h_1&=\frac{p_1}{b_0},\\
 h_2&=\frac1{b_0}\left(p_2-b_1h_1-\frac12b_0'h_1^2\right).
\end{align*}
They display both sources of nonlinear correction: the tail of the target
time form and the change of its ordinary leading coefficient along the
displacement.  The general recursion is the same calculation with a
well-ordered support rather than integer powers.

The support certificate is part of the result.  It tells us why each
coefficient in \eqref{eq:explicit-recursion} is defined even if the
multiples of the smallest input exponent are not cofinal.  It is not a
promise that arbitrary Hahn-series equality, primitive existence, or
period evaluation can be decided by a finite computer program.

\section{A compact hypothesis audit}
\label{app:audit}

\begin{longtable}{@{}p{0.28\textwidth}p{0.66\textwidth}@{}}
\toprule
Statement & Indispensable hypotheses and conclusion boundaries\\
\midrule
\endhead
Logarithm and exponential &
Set-sized ordered value group; common-domain holomorphic coefficients;
positive perturbative support.  No ordinary parameter convergence is
asserted.\\[5pt]
Difference equation &
$F\ne\id$.  If $W_F$ has a singular leading coefficient, the quotient
$g/W_F$ must still belong to the holomorphic Hahn algebra; period vanishing
alone is insufficient.\\[5pt]
Global conjugacy &
Same leading data $(\delta,a)$, $a$ nowhere zero, and a near-identity
conjugator on the original common domain.  Periods are ordinary
coefficientwise periods.\\[5pt]
Finite moduli &
The ordinary domain is the plane minus finitely many named points.
The basis consists of the specified puncture loops; arbitrary changes of
ordinary coordinates are not quotiented out.\\[5pt]
Flat nonconjugacy &
The chosen $\delta$ is noncofinal.  Jets are formed in
$\A_{\geq0}$; the ideal interpretation would be wrong in the full
algebra where $\eps$ is invertible.\\[5pt]
Centralizer &
A connected ordinary domain and a nonidentity map.  Permitted constants
satisfy $\val(cW_F)>0$; unrestricted large-time maps are not asserted.\\[5pt]
Slow time &
A specified ordinary holomorphic flow chart with image in $D$ and nonzero
variational factor $J$.  All Hahn coefficients exist on that chart;
continuation past its ordinary singularities is not supplied.\\
\bottomrule
\end{longtable}

\clearpage
\phantomsection
\addcontentsline{toc}{section}{References}
\begingroup
\small\raggedright
\begin{thebibliography}{9}

\bibitem{bkkps}
V. Bagayoko, L. S. Krapp, S. Kuhlmann, D. Panazzolo, and M. Serra,
\emph{Automorphisms and derivations on algebras endowed with formal infinite
sums}, arXiv:2403.05827, version 2, 2025 (first posted 2024).
\url{https://arxiv.org/abs/2403.05827}.

\bibitem{ab}
M. Aschenbrenner and W. Bergweiler,
\emph{Julia's equation and differential transcendence},
Illinois Journal of Mathematics \textbf{59} (2015), 277--294.
Preprint: \url{https://arxiv.org/abs/1307.6381}.

\bibitem{hairer}
E. Hairer, \emph{Geometric Numerical Integration, Lecture 3: Backward error
analysis}, TU M\"unchen lectures, January--February 2010.
\url{https://www.unige.ch/~hairer/poly_geoint/week3.pdf}.

\bibitem{gonshor}
H. Gonshor, \emph{An Introduction to the Theory of Surreal Numbers},
London Mathematical Society Lecture Note Series 110,
Cambridge University Press, 1986.
\url{https://www.cambridge.org/core/books/an-introduction-to-the-theory-of-surreal-numbers/312AE504A3E88E804054BFB390446374}.

\bibitem{repo-catalogue}
V. Reshetnikov, \emph{Surreal repository: research-document catalogue},
\texttt{docs/README.md}, snapshot
\texttt{39f2be6667ade51bca2b45daa47e289d69c09764}, consulted 21 September 2026.
\url{https://github.com/VladimirReshetnikov/Surreal/blob/39f2be6667ade51bca2b45daa47e289d69c09764/docs/README.md}.

\bibitem{repo-analysis}
\emph{Surcomplex analysis: holomorphic functions on $K=\No[i]$},
research-package documentation in the Surreal repository,
\texttt{docs/surcomplex/analysis/README.md}, same snapshot.
\url{https://github.com/VladimirReshetnikov/Surreal/blob/39f2be6667ade51bca2b45daa47e289d69c09764/docs/surcomplex/analysis/README.md}.

\bibitem{repo-contours}
\emph{Contours and Stokes: Jordan separation, Cauchy theory, and Stokes
calculus on the surcomplex plane}, research-package documentation in the
Surreal repository, \texttt{docs/surcomplex/contours-and-stokes/README.md},
same snapshot.
\url{https://github.com/VladimirReshetnikov/Surreal/blob/39f2be6667ade51bca2b45daa47e289d69c09764/docs/surcomplex/contours-and-stokes/README.md}.

\end{thebibliography}
\endgroup
\end{document}
